\chapter{Peano Arithmetic and Formal Number Theory}\label{ch:peano}

Having a sound and complete proof system for first-order logic, we can now pick a set of
axioms and study what follows. The most important example is arithmetic on the natural
numbers. This chapter also revisits induction, the main proof technique for reasoning
about programs and data structures.

\section{The language of arithmetic}

The language $\mathcal{L}_A$ has
\begin{itemize}
  \item a constant symbol $0$;
  \item a unary function symbol $s$ (successor; think $s(x)=x+1$);
  \item binary function symbols $+$ and $\cdot$ (written infix);
  \item equality $=$ (and sometimes a binary predicate $<$, which can also be defined).
\end{itemize}
The \emph{standard interpretation} $\mathcal{N}$ has domain $\N=\{0,1,2,\dots\}$ with the
usual meanings. The closed term $s(s(s(0)))$ denotes $3$; we abbreviate $s^n(0)$ by
$\overline{n}$ (a \emph{numeral}).

\section{The Peano axioms}

\begin{definition}[Peano arithmetic, PA]
\begin{align*}
&\text{(PA1)}\quad \forall x\,\neg(s(x)=0)\\
&\text{(PA2)}\quad \forall x\,\forall y\,(s(x)=s(y)\to x=y)\\
&\text{(PA3)}\quad \forall x\,(x+0=x)\\
&\text{(PA4)}\quad \forall x\,\forall y\,(x+s(y)=s(x+y))\\
&\text{(PA5)}\quad \forall x\,(x\cdot 0=0)\\
&\text{(PA6)}\quad \forall x\,\forall y\,(x\cdot s(y)=x\cdot y+x)\\
&\text{(PA7, induction scheme)}\quad
  \big(\varphi[0/x]\land\forall x\,(\varphi\to\varphi[s(x)/x])\big)\to\forall x\,\varphi
\end{align*}
for every formula $\varphi$ (which may have other free variables, implicitly universally
quantified).
\end{definition}

\begin{itemize}
  \item PA1 and PA2 say $0$ is not a successor and $s$ is injective, so
        $0, s(0), s(s(0)),\dots$ are all distinct.
  \item PA3--PA6 are the recursive definitions of $+$ and $\cdot$ (recursion on the second
        argument).
  \item PA7 is not one axiom but infinitely many: one for each formula $\varphi$. Its
        job is to rule out elements that are not reachable from $0$ by finitely many
        successor steps (as far as first-order logic can).
\end{itemize}

Order can be defined: $x\le y$ abbreviates $\exists z\,(x+z=y)$ and $x<y$ abbreviates
$\exists z\,(x+s(z)=y)$.

\section{Computing with the axioms}

\begin{example}[$\overline{2}+\overline{2}=\overline{4}$]
Informally, using PA4 twice and PA3 once:
\begin{align*}
s(s(0))+s(s(0)) &= s(s(s(0))+s(0)) && \text{PA4}\\
 &= s(s(s(s(0))+0)) && \text{PA4}\\
 &= s(s(s(s(0)))) && \text{PA3 and $=e$.}
\end{align*}
Every true equation between closed terms is provable this way: PA (in fact, the much
weaker system without induction) proves all true closed atomic sentences.
\end{example}

\section{Proofs by induction in PA}

To prove $\forall x\,\varphi$ in PA, prove the base case $\varphi[0/x]$ and the inductive
step $\forall x\,(\varphi\to\varphi[s(x)/x])$, then apply PA7 and $\to e$.

\begin{example}[$\vdash_{PA}\forall x\,(0+x=x)$]
Note this is \emph{not} an axiom: PA3 says $x+0=x$, with $0$ on the right. Let
$\varphi$ be $0+x=x$.
\begin{ndproof}
1 & $\forall x\,(x+0=x)$ & PA3\\
2 & $\forall x\,\forall y\,(x+s(y)=s(x+y))$ & PA4\\
3 & $0+0=0$ & $\forall e$ 1\\
4 & \bx\bx $0+x_0=x_0$ & $x_0$ fresh;\ assumption\\
5 & \bx\bx $\forall y\,(0+s(y)=s(0+y))$ & $\forall e$ 2\\
6 & \bx\bx $0+s(x_0)=s(0+x_0)$ & $\forall e$ 5\\
7 & \bx\bx $0+s(x_0)=s(x_0)$ & $=e$ 4, 6\\
8 & \bx $0+x_0=x_0\to 0+s(x_0)=s(x_0)$ & $\to i$ 4--7\\
9 & $\forall x\,(0+x=x\to 0+s(x)=s(x))$ & $\forall i$ 4--8\\
10 & $0+0=0\land\forall x\,(0+x=x\to 0+s(x)=s(x))$ & $\land i$ 3, 9\\
11 & $(0+0=0\land\forall x\,(\dots))\to\forall x\,(0+x=x)$ & PA7\\
12 & $\forall x\,(0+x=x)$ & $\to e$ 10, 11
\end{ndproof}
In line 7, $=e$ uses $\varphi'$ being $0+s(x_0)=s(z)$: line 6 is $\varphi'[0+x_0/z]$ and
line 4 says $0+x_0=x_0$, giving $\varphi'[x_0/z]$.
\end{example}

\begin{example}[$\forall x\,\neg(s(x)=x)$, informally]
Let $\varphi(x)$ be $\neg(s(x)=x)$.
\emph{Base}: $\neg(s(0)=0)$ is an instance of PA1.
\emph{Step}: assume $\neg(s(x_0)=x_0)$. If $s(s(x_0))=s(x_0)$, then by PA2,
$s(x_0)=x_0$, contradicting the IH. So $\neg(s(s(x_0))=s(x_0))$.
By PA7, $\forall x\,\neg(s(x)=x)$.
\end{example}

\begin{example}[Commutativity of $+$, sketch]
\begin{enumerate}
  \item Lemma A: $\forall x\,(0+x=x)$ (above).
  \item Lemma B: $\forall x\,\forall y\,(s(y)+x=s(y+x))$, by induction on $x$:
        base $s(y)+0=s(y)=s(y+0)$ by PA3; step
        $s(y)+s(x)=s(s(y)+x)=s(s(y+x))=s(y+s(x))$ by PA4, IH, PA4.
  \item Theorem: $\forall x\,\forall y\,(x+y=y+x)$, by induction on $y$:
        base $x+0=x=0+x$ by PA3 and Lemma A; step
        $x+s(y)=s(x+y)=s(y+x)=s(y)+x$ by PA4, IH, Lemma B.
\end{enumerate}
Choosing which variable to induct on, and which auxiliary lemmas to prove first, is the
creative part.
\end{example}

\section{Forms of induction}

\begin{important}[Equivalent induction principles on $\N$]
\begin{description}
  \item[Simple induction:] $P(0)$ and $\forall n\,(P(n)\to P(n+1))$ imply $\forall n\,P(n)$.
  \item[Strong (complete) induction:] if $\forall n\,\big((\forall k<n\ P(k))\to P(n)\big)$
        then $\forall n\,P(n)$. (No separate base case is needed, but the case $n=0$ must
        be handled with no hypotheses.)
  \item[Well-ordering:] every non-empty subset of $\N$ has a least element. A proof by
        ``minimal counterexample'' uses this.
\end{description}
Strong induction follows from simple induction applied to $Q(n)=\forall k\le n\ P(k)$.
\end{important}

Structural induction (Chapter~\ref{ch:propsyntax}) generalises these to any inductively
defined set: formulas, lists, trees. Induction on the length of a proof (used in the
soundness proof) is strong induction on $\N$.

\begin{example}[Strong induction]
Every integer $n\ge2$ is a product of primes. If $n$ is prime, done. Otherwise $n=ab$ with
$2\le a,b<n$; by IH both $a$ and $b$ are products of primes, hence so is $n$.
\end{example}

\begin{important}[Common pitfalls]
\begin{itemize}
  \item Forgetting the base case (``$n=n+1$ implies $n+1=n+2$'' proves nothing).
  \item An inductive step that silently fails for small $n$ (the classic ``all horses are
        the same colour'' argument breaks going from $1$ to $2$).
  \item Using the conclusion $\forall x\,\varphi$ inside the step: the IH is only
        $\varphi$ for the particular $x_0$.
  \item Inducting on the wrong variable, or with a statement too weak to carry the step;
        strengthening the claim (``loading the induction'') often helps.
\end{itemize}
\end{important}

\section{Models of PA and the limits of first-order arithmetic}

\begin{itemize}
  \item The standard model $\mathcal{N}$ satisfies all axioms of PA, so (by soundness)
        PA is consistent if we accept that $\mathcal{N}$ exists.
  \item \textbf{Non-standard models.} Add a new constant $c$ and sentences
        $\overline{n}<c$ for every $n$. Every finite subset is satisfied in $\mathcal{N}$
        (interpret $c$ as a large number), so by compactness the whole set has a model:
        a model of PA with an element larger than every numeral. So first-order PA cannot
        pin down $\N$ up to isomorphism. (Second-order induction, quantifying over all
        subsets, can.)
  \item \textbf{Incompleteness.} G\"odel (1931) showed there are sentences true in
        $\mathcal{N}$ that PA cannot prove, and that this cannot be fixed by adding more
        (computably listable) axioms. See Chapter~\ref{ch:undecidability}.
  \item By contrast, \emph{Presburger arithmetic} (only $0$, $s$, $+$, no $\cdot$) is
        complete and decidable.
\end{itemize}

\begin{aside}
The connection to computer science: the recursive equations PA3--PA6 are exactly how
addition and multiplication would be programmed on unary numbers:
\begin{lstlisting}[style=racket]
(define (add x y)            ; recursion on y, like PA3/PA4
  (cond [(zero? y) x]
        [else (add1 (add x (sub1 y)))]))
\end{lstlisting}
and an induction proof of a property of \code{add} follows the shape of its recursion.
\end{aside}

\section*{Summary}
\begin{itemize}
  \item PA: $0$ is not a successor; $s$ is injective; recursive definitions of $+$ and
        $\cdot$; and the induction scheme (one axiom per formula).
  \item Facts like $0+x=x$ and $x+y=y+x$ need induction; closed equations just need
        PA3--PA6.
  \item Simple induction, strong induction and well-ordering are equivalent on $\N$;
        structural induction generalises them.
  \item First-order PA has non-standard models (compactness) and is incomplete (G\"odel).
\end{itemize}
